\section{Characteristic functions, Fourier formulas, and applications}\label{sec:fourier}

\subsection{Characteristic functions and the first moment}

For a finite rational law, define the regular complex name
\[
 \phi_X(t)=\E[e^{itX}]
 =\E[\cos(tX)]+i\E[\sin(tX)].
\]
The sine and cosine names are obtained from finite Taylor sums with explicit
remainders after range reduction.

\begin{theorem}[Derivative of the characteristic function at zero]
\label{thm:char-derivative}
If \(\E[X^2]\) is finite with an effective name, then for every nonzero
rational \(h\),
\[
 \left|\frac{\phi_X(h)-1}{h}-i\E[X]\right|
 \le\frac{\abs h}{2}\E[X^2].
\]
Consequently \(\phi_X'(0)=i\E[X]\) effectively.
\end{theorem}

\begin{proof}
The elementary Taylor remainder estimate
\[
 \abs{e^{iu}-1-iu}\le\frac{u^2}{2}
\]
is proved constructively from the integral or alternating-series remainder for
sine and cosine; at rational precision it is a finite polynomial inequality
plus a geometric tail bound.  Put \(u=hX_i\), divide by \(\abs h\), and take
expectations:
\begin{align*}
\left|\frac{\phi_X(h)-1}{h}-i\E[X]\right|
&=\left|\E\left[\frac{e^{ihX}-1-ihX}{h}\right]\right|\\
&\le\E\left[\frac{\abs{e^{ihX}-1-ihX}}{\abs h}\right]\\
&\le\frac{\abs h}{2}\E[X^2].
\end{align*}
For an effective infinite law, first truncate using the second-moment tail
modulus.  Given output precision, choose \(h\) small enough that the final
bound is below it.
\end{proof}

\subsection{Finite cyclic Fourier inversion}

\begin{theorem}[Finite Fourier expectation identity]\label{thm:finite-fourier}
Let \(p_x\) be a rational probability mass on \(\Z/N\Z\), let
\(\omega=e^{2\pi i/N}\), and define
\[
 \phi(k)=\sum_{x<N}p_x\omega^{kx},\qquad
 \widehat g(k)=\sum_{x<N}g(x)\omega^{-kx}.
\]
Then
\[
 \E[g(X)]=\frac1N\sum_{k<N}\widehat g(k)\phi(k).
\]
This identity has an effective \(\ipra\) realization using regular complex
names, and an exact algebraic realization in a cyclotomic presentation.
\end{theorem}

\begin{proof}
Expand the right side and use finite Fubini:
\[
 \frac1N\sum_{x,y<N}g(x)p_y
       \sum_{k<N}\omega^{k(y-x)}.
\]
The inner sum is \(N\) if \(x=y\) and zero otherwise.  For \(x\ne y\), this
is the finite geometric identity
\[
 (1-\omega^{y-x})\sum_{k<N}\omega^{k(y-x)}
 =1-\omega^{N(y-x)}=0,
\]
together with a primitive-root certificate that
\(1-\omega^{y-x}\ne0\).  Thus only diagonal terms remain, giving
\(\sum_xg(x)p_x\).  With regular complex names, approximate all roots of unity
with guard bits and use the same exact exponent arithmetic; with a cyclotomic
algebra code, the orthogonality certificate is symbolic.
\end{proof}

\subsection{Effective continuous Fourier formula}

\begin{theorem}[Expectation through Fourier inversion]
\label{thm:continuous-fourier}
Suppose a density \(f\), its characteristic function \(\phi\), a function
\(g\), and its Fourier transform
\[
 G(t)=\int_{\R}g(x)e^{-itx}\,dx
\]
are given by effective integration names.  Assume explicit absolute-tail and
Fubini moduli for the double integral.  Then, with the article's sign
convention,
\[
 \E[g(X)]=\frac1{2\pi}\int_{\R}G(t)\phi(t)\,dt.
\]
\end{theorem}

\begin{proof}
Truncate both \(x\) and \(t\) to compact rational intervals chosen from the
absolute-tail modulus.  Approximate both integrals by rational rectangles.
On the resulting finite double sum, interchange the order by exact finite
Fubini and substitute the finite Fourier inversion approximation to \(f\).
The accumulated quadrature, inversion, and tail errors are allocated less than
\(2^{-p}/4\) each.  Since the moduli compute all truncation radii and mesh
sizes from \(p\), this is an effective theorem schema.  Without absolute
Fubini data, the order change is not claimed.
\end{proof}

\subsection{Moments and finite moment determinacy}

\begin{definition}[Moments]
The \(r\)-th raw moment is \(m_r=\E[X^r]\), and the \(r\)-th central moment is
\(\E[(X-\E X)^r]\).  For finite rational laws these are exact rational
expectations.  Their finite lists are primitive recursive in the law code.
\end{definition}

\begin{theorem}[Finite-support moment determinacy]\label{thm:moment-determinacy}
Let \(X\) be supported on \(\{0,1,\ldots,B\}\).  The moments
\(m_0,m_1,\ldots,m_B\) determine the probability masses uniquely, and the
masses are primitive-recursively recoverable by rational arithmetic.
\end{theorem}

\begin{proof}
Use falling factorials
\((x)_r=x(x-1)\cdots(x-r+1)\) and factorial moments
\(f_r=\E[(X)_r]\).  Ordinary powers and falling factorials are related by the
finite triangular Stirling-number transformations, whose integer coefficients
are primitive recursive.  Hence \((m_0,\ldots,m_B)\) determines
\((f_0,\ldots,f_B)\).

If \(p_k=\Pp(X=k)\), then
\[
 f_r=\sum_{k=r}^{B}(k)_r p_k.
\]
This is upper triangular, with diagonal \((r)_r=r!>0\).  Recover the masses by
back substitution:
\[
 p_B=\frac{f_B}{B!},\qquad
 p_r=\frac{f_r-\sum_{k=r+1}^{B}(k)_r p_k}{r!}.
\]
A downward bounded recursion computes every \(p_r\).  Uniqueness follows from
the same recursion applied to two candidate vectors.  This is a complete
\(\ipra\) proof of moment determinacy for known finite support; no analytic
moment problem is involved.
\end{proof}

\begin{corollary}[Finite moment generating function]
For finite support, the formal polynomial/exponential sum
\(M(t)=\E[e^{tX}]\) determines the law, since its first \(B\) derivatives at
zero give the moments in \Cref{thm:moment-determinacy}.
\end{corollary}

\subsection{Expected utility and finite decision problems}

\begin{theorem}[Finite expected-utility maximization]\label{thm:expected-utility}
Let \(A\) be a finite set of actions, \(S\) a finite weighted state law, and
\(u(a,s)\in\Q\) a primitive-recursive utility table.  There is a
primitive-recursively selected action \(a^*\) satisfying
\[
 \E[u(a^*,S)]\ge\E[u(a,S)]\quad\text{for every }a\in A.
\]
\end{theorem}

\begin{proof}
Compute the exact rational expected utility for every action.  Bounded search
selects the least index attaining the maximum, using decidable rational order.
An induction over the action list proves that the running candidate dominates
all earlier candidates.  The output includes its comparison certificate.
\end{proof}

\begin{remark}
This proves the computational content of finite expected-utility choice.  It
does not assert that every preference ordering has a utility representation;
that is a separate representation theorem with its own hypotheses.
\end{remark}

\subsection{Indicators, Monte Carlo, and empirical means}

\begin{theorem}[Probability as expectation and unbiased frequency]
\label{thm:monte-carlo}
For an event \(A\), \(\Pp(A)=\E[\ind_A]\).  For iid repetitions and empirical
frequency
\[
 \widehat p_n=\frac1n\sum_{j<n}\ind_A^{(j)},
\]
we have
\[
 \E[\widehat p_n]=\Pp(A),\qquad
 \Var(\widehat p_n)=\frac{p(1-p)}n.
\]
Consequently Chebyshev supplies the exact finite Monte Carlo error contract
\[
 \Pp(\abs{\widehat p_n-p}\ge\varepsilon)
 \le\frac{p(1-p)}{n\varepsilon^2}
 \le\frac1{4n\varepsilon^2}.
\]
\end{theorem}

\begin{proof}
The indicator identity was proved in \Cref{thm:indicator}.  Apply the
sample-mean and variance theorems to Bernoulli variables, whose variance is
\(p-p^2=p(1-p)\).  Finally,
\(p(1-p)\le1/4\) is the square identity
\(1/4-p(1-p)=(p-1/2)^2\ge0\), interpreted after clearing denominators.
\end{proof}

\subsection{Center of mass}

\begin{theorem}[Balance point]\label{thm:center-mass}
Place nonnegative rational masses \(m_i\) at rational positions \(x_i\), with
positive total mass \(M=\sum_i m_i\).  The unique point \(c\) with zero total
torque
\[
 \sum_i m_i(x_i-c)=0
\]
is
\[
 c=\frac{\sum_i m_ix_i}{M}.
\]
Thus the center of mass is exactly the expectation under probabilities
\(p_i=m_i/M\).
\end{theorem}

\begin{proof}
Expand the torque:
\[
 \sum_i m_ix_i-c\sum_i m_i=\sum_i m_ix_i-cM.
\]
It vanishes at the displayed \(c\).  Conversely, if it vanishes at \(d\), then
\(dM=\sum_i m_ix_i=cM\); positive \(M\) permits cancellation.  This is a
single denominator-cleared linear identity.
\end{proof}

\begin{proposition}[Beta balance example]\label{prop:beta}
For positive natural \(\alpha,\beta\), the beta density proportional to
\(x^{\alpha-1}(1-x)^{\beta-1}\) on \([0,1]\) has mean
\(\alpha/(\alpha+\beta)\).
\end{proposition}

\begin{proof}
For natural parameters, the density is a polynomial times a rational
normalization.  Integration by parts gives the finite beta recursion
\[
 B(\alpha+1,\beta)=\frac{\alpha}{\alpha+\beta}B(\alpha,\beta).
\]
The expectation is the ratio of these two exact polynomial integrals.  The
recursion can alternatively be proved by expanding
\((1-x)^{\beta-1}\) and integrating monomials rationally.
\end{proof}

\subsection{A bounded Wald identity}

\begin{theorem}[Wald equation for a bounded stopping time]
\label{thm:wald}
Let \(X_0,\ldots,X_{N-1}\) be iid finite rational variables with mean \(\mu\).
Let \(T\le N\) be a stopping time, meaning the event \(T>j\) depends only on
\(X_0,\ldots,X_{j-1}\).  Then
\[
 \E\left[\sum_{j<T}X_j\right]=\mu\E[T].
\]
\end{theorem}

\begin{proof}
Pointwise,
\[
 \sum_{j<T}X_j=\sum_{j<N}X_j\ind_{T>j},
 \qquad
 T=\sum_{j<N}\ind_{T>j}.
\]
Finite linearity reduces the first expectation to a sum.  By the stopping-time
condition, \(\ind_{T>j}\) is a function of earlier coordinates and is
independent of \(X_j\).  Therefore
\[
 \E[X_j\ind_{T>j}]=\E[X_j]\Pp(T>j)=\mu\Pp(T>j).
\]
Sum over \(j<N\) and use the tail formula for the bounded integer \(T\).
\end{proof}

\subsection{Finite-dimensional quantum expectation}

\begin{theorem}[Quantum expectation and variance identity]
\label{thm:quantum}
Let \(\psi\in\Q(i)^d\) satisfy \(\psi^*\psi=1\), and let
\(A\in\Q(i)^{d\times d}\) be Hermitian.  Put
\[
 \langle A\rangle=\psi^*A\psi.
\]
Then \(\langle A\rangle\) is real and
\[
 (\Delta A)^2
 \defeq\psi^*A^2\psi-(\psi^*A\psi)^2
 =\norm{(A-\langle A\rangle I)\psi}^2\ge0.
\]
\end{theorem}

\begin{proof}
Hermiticity gives
\((\psi^*A\psi)^*=\psi^*A^*\psi=\psi^*A\psi\), so the scalar is real.  Expand
the squared norm:
\begin{align*}
\psi^*(A-\mu I)^*(A-\mu I)\psi
 &=\psi^*(A^2-2\mu A+\mu^2I)\psi\\
 &=\psi^*A^2\psi-2\mu^2+\mu^2
 =\psi^*A^2\psi-\mu^2,
\end{align*}
where \(\mu=\psi^*A\psi\) and \(\psi^*\psi=1\).  The norm square is a finite
sum of rational squares of real and imaginary parts, hence nonnegative in
\(\ipra\).  Algebraic or regular complex entries are handled by the same
finite identity at each precision.
\end{proof}
